\chapter{Результаты эксперимента} \label{ch3}

\section{Условия эксперимента} \label{ch3:conditions}

Эксперимент выполнялся в локальном каталоге \texttt{benchmark\_app}. Для каждого формата входа выполнялось по три запуска каждого фреймворка. В таблицах ниже приведены медианные значения по трем повторам. Все процессы завершились с кодом \texttt{0}, то есть тесты Jest, Vitest и Mocha прошли успешно.

Следует учитывать, что проценты CPU в \texttt{pidusage} могут превышать 100\%, если процесс использует несколько логических ядер. Поэтому значение CPU интерпретируется как пиковая интенсивность использования процессора, а не как доля одного ядра.

\section{Результаты для JSON-входа} \label{ch3:json}

JSON-набор \texttt{author-json-medium} имеет размер \texttt{5000}. Медианные результаты приведены в \taref{tab:json-results}.

\noindent
\begingroup
\centering
\small
\begin{longtable}[c]{|p{3.5cm}|p{3.2cm}|p{3.2cm}|p{3.2cm}|}
    \caption{Медианные результаты для JSON-входа}
    \label{tab:json-results}
    \\
    \hline
    Фреймворк & Время, мс & Пиковая память, МБ & Пиковый CPU, \% \\ \hline
    \endfirsthead
    \hline
    Фреймворк & Время, мс & Пиковая память, МБ & Пиковый CPU, \% \\ \hline
    \endhead
    \hline
    Jest & 10640.290 & 90.781 & 36.000 \\ \hline
    Vitest & 16273.677 & 372.242 & 136.860 \\ \hline
    Mocha & 3606.863 & 57.340 & 4.700 \\ \hline
\end{longtable}
\normalsize
\endgroup

На данном наборе Mocha оказался самым быстрым и самым легким по памяти. Vitest показал наибольшую пиковую память, что связано с инфраструктурой Vite/Vitest и пулом выполнения. Jest занял промежуточное положение по памяти и времени.

\section{Результаты для CSV-входа} \label{ch3:csv}

CSV-набор \texttt{author-example.csv} содержит 15 значений. Медианные результаты приведены в \taref{tab:csv-results}.

\noindent
\begingroup
\centering
\small
\begin{longtable}[c]{|p{3.5cm}|p{3.2cm}|p{3.2cm}|p{3.2cm}|}
    \caption{Медианные результаты для CSV-входа}
    \label{tab:csv-results}
    \\
    \hline
    Фреймворк & Время, мс & Пиковая память, МБ & Пиковый CPU, \% \\ \hline
    \endfirsthead
    \hline
    Фреймворк & Время, мс & Пиковая память, МБ & Пиковый CPU, \% \\ \hline
    \endhead
    \hline
    Jest & 10754.265 & 87.238 & 26.700 \\ \hline
    Vitest & 14242.386 & 370.348 & 326.500 \\ \hline
    Mocha & 3559.561 & 57.242 & 0.000 \\ \hline
\end{longtable}
\normalsize
\endgroup

Для малого CSV-набора сохраняется та же общая картина: Mocha имеет минимальное время и память, Jest находится между Mocha и Vitest, Vitest использует больше памяти. Нулевое медианное значение CPU у Mocha означает, что два из трех запусков завершились между моментами сэмплирования CPU. Это ограничение дискретного мониторинга, а не доказательство отсутствия процессорной нагрузки.

\section{Сводное сравнение} \label{ch3:summary}

\noindent
\begingroup
\centering
\small
\begin{longtable}[c]{|p{3.2cm}|p{3.2cm}|p{3.2cm}|p{4.2cm}|}
    \caption{Лучшие значения по метрикам}
    \label{tab:best-values}
    \\
    \hline
    Набор & Лучшая скорость & Лучшая память & Комментарий \\ \hline
    \endfirsthead
    \hline
    Набор & Лучшая скорость & Лучшая память & Комментарий \\ \hline
    \endhead
    \hline
    JSON, 5000 элементов & Mocha & Mocha & Минимальная инфраструктура дала преимущество на малом тестовом наборе. \\ \hline
    CSV, 15 элементов & Mocha & Mocha & Для короткого теста стоимость запуска раннера доминирует над вычислениями. \\ \hline
\end{longtable}
\normalsize
\endgroup

В рамках данного эксперимента Mocha оказался наиболее экономичным. Однако это не означает, что Mocha всегда является лучшим выбором. В реальном проекте важны также mocks, coverage, watch-режим, интеграция со сборщиком, привычность команды и существующие тесты.
